\documentclass[12pt,a4paper]{article}
\usepackage[utf8]{inputenc}
\usepackage[indonesian]{babel}
\usepackage{amsmath}
\usepackage{amsfonts}
\usepackage{amssymb}
\usepackage{geometry}

\geometry{margin=2.5cm}

\title{Menghitung Luas Permukaan Benda Putar Sumbu-x}
\author{Ahmad Fakhrul Bawani}
\date{}

\begin{document}

\maketitle

\section*{Soal}

Find the area of the surface generated by revolving the curve:
\begin{equation*}
  y = \sqrt{x+5}, \quad \text{untuk } 1 \le x \le 7, \quad \text{mengelilingi sumbu-}x
\end{equation*}

\subsection*{1. Rumus Umum}
Rumus luas permukaan ($S$) benda putar yang dihasilkan oleh perputaran mengelilingi sumbu-$x$ dari $x = a$ sampai $x = b$ didefinisikan sebagai:
\begin{equation*}
  S = 2\pi \int_{a}^{b} y \cdot \sqrt{1 + \left(\frac{dy}{dx}\right)^2} \, dx
\end{equation*}

\subsection*{2. Menentukan Turunan Fungsi}
Tulis ulang fungsi ke dalam bentuk pangkat: $y = (x+5)^{\frac{1}{2}}$. Cari turunan pertama $y$ terhadap $x$:
\begin{equation*}
  \frac{dy}{dx} = \frac{1}{2}(x+5)^{-\frac{1}{2}} = \frac{1}{2\sqrt{x+5}}
\end{equation*}

Kuadratkan hasil turunan tersebut:
\begin{equation*}
  \left(\frac{dy}{dx}\right)^2 = \left(\frac{1}{2\sqrt{x+5}}\right)^2 = \frac{1}{4(x+5)} = \frac{1}{4x+20}
\end{equation*}

Masukkan ke dalam komponen akar panjang busur dan samakan penyebutnya:
\begin{equation*}
  \sqrt{1 + \left(\frac{dy}{dx}\right)^2} = \sqrt{1 + \frac{1}{4(x+5)}} = \sqrt{\frac{4(x+5) + 1}{4(x+5)}} = \sqrt{\frac{4x+21}{4(x+5)}} = \frac{\sqrt{4x+21}}{2\sqrt{x+5}}
\end{equation*}

\subsection*{3. Solusi Integrasi}
Substitusikan komponen $y = \sqrt{x+5}$ dan bentuk akar di atas ke dalam rumus luas permukaan dengan batas $a = 1$ dan $b = 7$:
\begin{align*}
  S &= 2\pi \int_{1}^{7} \left(\sqrt{x+5}\right) \cdot \left(\frac{\sqrt{4x+21}}{2\sqrt{x+5}}\right) \, dx
\end{align*}
Suku $\sqrt{x+5}$ pada pembilang dan penyebut saling menghilangkan, dan angka 2 pada penyebut membagi habis konstanta $2\pi$:
\begin{align*}
  S &= \pi \int_{1}^{7} \sqrt{4x+21} \, dx \\
  S &= \pi \int_{1}^{7} (4x+21)^{\frac{1}{2}} \, dx
\end{align*}

Gunakan metode substitusi untuk menyelesaikan integral tersebut:
\begin{align*}
  \text{Misalkan: } & u = 4x + 21 \implies du = 4 \, dx \implies dx = \frac{1}{4} \, du \\
  \text{Batas baru: } & \text{Untuk } x = 1 \implies u = 4(1) + 21 = 25 \\
  & \text{Untuk } x = 7 \implies u = 4(7) + 21 = 28 + 21 = 49
\end{align*}

Substitusikan variabel baru $u$ ke dalam integral:
\begin{align*}
  S &= \pi \int_{25}^{49} u^{\frac{1}{2}} \cdot \left(\frac{1}{4} \, du\right) \\
  S &= \frac{\pi}{4} \int_{25}^{49} u^{\frac{1}{2}} \, du \\
  S &= \frac{\pi}{4} \left[ \frac{2}{3}u^{\frac{3}{2}} \right]_{25}^{49} \\
  S &= \frac{\pi}{6} \left[ u\sqrt{u} \right]_{25}^{49} \\
  S &= \frac{\pi}{6} \left( \left[ 49\sqrt{49} \right] - \left[ 25\sqrt{25} \right] \right) \\
  S &= \frac{\pi}{6} \left( 49(7) - 25(5) \right) \\
  S &= \frac{\pi}{6} (343 - 125) \\
  S &= \frac{\pi}{6} (218) \\
  S &= \frac{109\pi}{3}
\end{align*}

Jadi, untuk luas permukaan benda putar tersebut adalah \textbf{$\dfrac{109\pi}{3}$ satuan luas}.

\end{document}